% Part: turing-machines
% Chapter: undecidability
% Section: halting-problem

\documentclass[../../../include/open-logic-section]{subfiles}

\begin{document}

\olfileid{tur}{und}{hal}
\olsection{અટકવાની સમસ્યા}

\begin{explain}
માનો કે આપણે ટ્યુરિંગ મશીનોનાં વર્ણનોની કોઈ પરિગણના
નક્કી કરી છે. તેથી દરેક ટ્યુરિંગ મશીનને એક
\emph{સૂચક} મળે છે: મશીનોનાં વર્ણનોની પરિગણના
$M_1$, $M_2$, $M_3$, \dots{}માં તેનું સ્થાન.

ટ્યુરિંગ-અસંગણનીય વિધેયો હોવાં જ જોઈએ એ આપણે
જાણીએ છીએ: ટ્યુરિંગ મશીનોનાં વર્ણનોનો ગણ---અને તેથી
ટ્યુરિંગ મશીનોનો ગણ---!!{enumerable} છે, પરંતુ
$\Nat$થી $\Nat$માં જતાં બધાં વિધેયોનો ગણ એવો નથી.
તેમ છતાં આપણે અસંગણનીય વિધેયોનાં ચોક્કસ ઉદાહરણો પણ
મેળવી શકીએ. આવું એક વિધેય અટકવાનું વિધેય છે.
\end{explain}

\begin{defn}[અટકવાનું વિધેય]
 \emph{અટકવાનું વિધેય}~$h$ આ રીતે વ્યાખ્યાયિત છે:
\[
h(e,n) =
\begin{cases}
  \text{0} & \text{જો મશીન~$M_e$ ઇનપુટ $n$ પર ન અટકે} \\
  \text{1} & \text{જો મશીન~$M_e$ ઇનપુટ $n$ પર અટકે}
\end{cases}
\]
\end{defn}

\begin{defn}[અટકવાની સમસ્યા]
\emph{અટકવાની સમસ્યા} એટલે કોઈ પણ $e$, $n$ માટે,
ટ્યુરિંગ મશીન~$M_e$~$n$ સ્ટ્રોકના
ઇનપુટથી શરૂ કરતાં અટકે છે કે નહીં તે નક્કી કરવાની
સમસ્યા.
\end{defn}

\begin{explain}
સંબંધિત વિધેય~$s$ ટ્યુરિંગ-સંગણનીય નથી એવું બતાવીને
આપણે~$h$ પણ ટ્યુરિંગ-સંગણનીય નથી એવું બતાવીએ છીએ.
આ સાબિતી એ હકીકતો પર આધાર રાખે છે કે ટ્યુરિંગ મશીન
વડે ગણી શકાય તે બધું શિસ્તબદ્ધ ટ્યુરિંગ મશીન વડે પણ
ગણી શકાય છે (\olref[mac][dis]{sec}), અને બે ટ્યુરિંગ
મશીનોને જોડીને એક મશીન બનાવી શકાય છે
(\olref[mac][cmb]{sec}).
\end{explain}

\begin{defn} વિધેય~$s$ આ રીતે વ્યાખ્યાયિત છે:
\[
s(e) =
\begin{cases}
  \text{0} & \text{જો મશીન~$M_e$ ઇનપુટ $e$ પર ન અટકે} \\
  \text{1} & \text{જો મશીન~$M_e$ ઇનપુટ $e$ પર અટકે}
\end{cases}
\]
\end{defn}

\begin{lem}
વિધેય~$s$ ટ્યુરિંગ-સંગણનીય નથી.
\end{lem}

\begin{proof}
વિરોધાભાસ માટે માનો કે વિધેય~$s$ ટ્યુરિંગ-સંગણનીય
છે. તો~$s$ ગણતું ટ્યુરિંગ મશીન~$S$ હોત. વ્યાપકતા
ગુમાવ્યા વિના ધારી શકીએ કે $S$ અટકે ત્યારે પ્રથમ
ખાનાને વાંચતું હોય (એટલે તે શિસ્તબદ્ધ હોય). આ
મશીનને બીજા મશીન~$J$ સાથે ``જોડી'' શકાય. તે ઇનપુટ
~$0$ પર શરૂ થાય તો અટકે છે (એટલે શરૂઆતની અવસ્થામાં,
ટેપના છેડાના ચિહ્નની જમણી બાજુના ખાનાને વાંચતી વખતે
$\TMblank$ વાંચે તો અટકે છે); અન્યથા તે કદી ન અટકીને
જમણી બાજુ ભટકતું રહે છે. $S$ને~$J$ સાથે જોડીને
બનેલું મશીન $S \concat J$ પણ ટ્યુરિંગ મશીન છે;
તેથી કોઈ~$e$ માટે તે $M_e$ છે (એટલે તે પરિગણનામાં
ક્યાંક આવે છે). હવે~$M_e$ને~$e$ જેટલા $\TMstroke$ના
ઇનપુટથી શરૂ કરો. બે શક્યતાઓ છે: $M_e$ અટકે છે
અથવા નથી અટકતું.
\begin{enumerate}
\item માનો કે~$M_e$,~$e$ જેટલા $\TMstroke$ના ઇનપુટ
  પર અટકે છે. તો $s(e) = 1$. તેથી~$S$ને ઇનપુટ~$e$
  પર શરૂ કરતાં તે ટેપ પર આઉટપુટ તરીકે એક જ
  $\TMstroke$ રાખીને અટકે છે. ત્યાર પછી~$J$ ટેપ પર
  $\TMstroke$ હોય ત્યારે શરૂ થાય છે. આ કિસ્સામાં~$J$
  અટકતું નથી. પરંતુ~$M_e$ એ $S \concat J$ મશીન જ
  છે, તેથી તેને $S$ અને પછી $J$ જે કરે તે જ કરવું
  પડે (એટલે આ કિસ્સામાં જમણી બાજુ ભટકીને કદી ન
  અટકવું પડે). તેથી~$e$ જેટલા $\TMstroke$ના ઇનપુટ
  પર~$M_e$ અટકી શકે નહીં.

\item હવે માનો કે~$M_e$,~$e$ જેટલા $\TMstroke$ના
  ઇનપુટ પર અટકતું નથી. તો $s(e) = 0$, અને~$S$ને
  ઇનપુટ~$e$ પર શરૂ કરતાં તે ખાલી ટેપ સાથે અટકે છે.
  ખાલી ટેપ પર શરૂ કરાયેલું~$J$ તરત અટકે છે. ફરી,
  $M_e$ એ $S$ અને પછી~$J$ જે કરે તે જ કરે છે;
  તેથી~$e$ જેટલા $\TMstroke$ના ઇનપુટ પર~$M_e$
  અટકવું જ જોઈએ.
\end{enumerate}
બંને કિસ્સામાં ધારણા સાથે વિરોધાભાસ મળે છે. તેથી
આવું ટ્યુરિંગ મશીન~$S$ હોઈ શકે નહીં: $s$ ટ્યુરિંગ-સંગણનીય
નથી.
\end{proof}

\begin{thm}[અટકવાની સમસ્યા ઉકેલી ન શકાય]
\ollabel{thm:halting-problem} અટકવાની સમસ્યા ઉકેલી
શકાતી નથી, એટલે વિધેય~$h$ ટ્યુરિંગ-સંગણનીય નથી.
\end{thm}

\begin{proof}
માનો કે~$h$ ટ્યુરિંગ-સંગણનીય છે અને તેને ટ્યુરિંગ
મશીન~$H$ ગણે છે. $H$ વડે~$s$ ગણતું ટ્યુરિંગ મશીન
બનાવી શકીએ: પહેલાં ઇનપુટની એક નકલ કરો (તેને
~$\TMblank$ ચિહ્નથી છૂટી રાખીને); પછી શરૂઆતમાં
પાછા જઈ~$H$ ચલાવો. પહેલું કામ કરતું મશીન સ્પષ્ટપણે
બનાવી શકાય છે (\cref{tur:mac:dis:prob:copier});
અને જો $H$ હોત, તો તેને આવા નકલ કરનારા મશીન સાથે
``જોડીને'' એવું નવું મશીન મેળવી શકાત જે~$M_e$
ઇનપુટ~$e$ પર અટકે છે કે નહીં તે નક્કી કરે, એટલે~$s$
ગણે. પરંતુ આવું કોઈ મશીન હોઈ શકે નહીં એ આપણે પહેલેથી
બતાવ્યું છે. તેથી~$h$ પણ ટ્યુરિંગ-સંગણનીય નથી.
\end{proof}

\begin{prob}
ત્રણ-અટક (3-Halt) સમસ્યા એટલે અન્યથા ખાલી ટેપ પર
ત્રણ $\TMstroke$ના ઇનપુટ માટે મનસ્વી પસંદ કરેલું
ટ્યુરિંગ મશીન અટકે છે કે નહીં તે નક્કી કરતી નિર્ણય-વિધિ
આપવાની સમસ્યા. ત્રણ-અટક સમસ્યા ઉકેલી શકાતી નથી
તે સાબિત કરો.
\end{prob}

\begin{prob}
બતાવો કે~$e \neq n$ હોય તેવા ટ્યુરિંગ મશીન અને ઇનપુટ
જોડ $M_e$ તથા~$n$ માટે અટકવાની સમસ્યા ઉકેલી શકાય,
તો $e = n$ હોય તેવા કિસ્સાઓ માટે પણ ઉકેલી શકાય છે.
\end{prob}

\begin{prob}
અમે સાબિત કર્યું કે ઇનપુટ તરીકે સંખ્યા~$e$ આપવામાં આવે,
જે બધાં ટ્યુરિંગ મશીનોની પરિગણનામાં મશીન~$M_e$ને
ઓળખાવે છે, ત્યારે અટકવાની સમસ્યા ઉકેલી શકાતી નથી.
જો ઇનપુટ તરીકે~\olref[tur][und][enu]{sec}માંનાં
ટ્યુરિંગ મશીનોનાં વર્ણનો સીધાં આપવાની છૂટ આપીએ તો?
શું સૂચકોને બદલે મશીનોનાં વર્ણનો ઇનપુટ તરીકે લેતું,
અટકવાની સમસ્યાનો નિર્ણય કરતું ટ્યુરિંગ મશીન હોઈ
શકે? કેમ અથવા કેમ નહીં તે સમજાવો.
\end{prob}

\begin{prob} બતાવો કે નીચે પ્રમાણે વ્યાખ્યાયિત
  \emph{આંશિક} વિધેય~$s'$ ટ્યુરિંગ-સંગણનીય \emph{છે}:
  \[
  s'(e) =
  \begin{cases}
    \text{1} & \text{જો મશીન~$M_e$ ઇનપુટ $e$ પર અટકે}\\
    \text{અવ્યાખ્યાયિત} & \text{જો મશીન~$M_e$ ઇનપુટ $e$ પર ન અટકે}
  \end{cases}
  \]
\end{prob}

\end{document}
